% 4. EXPERIMENTAL SETUP
% =============================================================================
\section{Experimental Setup}
\label{sec:experiments_setup}

\subsection{Datasets and Diagnostic Pathology Coverage}
\label{sec:datasets}
We evaluate BioDINO-Edge across two distinct foliar pathology regimes:

\subsubsection{PlantVillage Controlled Benchmark (Laboratory Regime)}
From the standard PlantVillage repository \cite{hughes2012plantvillage}, we build a curated 11-class subset of 9,385 foliar images from three crops: Apple (\textit{Malus domestica}), Grape (\textit{Vitis vinifera}), and Potato (\textit{Solanum tuberosum}). The classes cover healthy foliage and eight diseases: Apple Scab, Apple Black Rot, Apple Cedar Rust, Grape Black Rot, Grape Esca, Grape Leaf Blight, Potato Early Blight, and Potato Late Blight. The subset is chosen from the 38 available classes for a specific reason: each host contributes at least two visually similar foliar pathologies plus a healthy class, so the task tests within-species confusion rather than coarse species separation. Any method that separates crops without resolving lesions will therefore score poorly, and the false-negative rate is measured against a genuine healthy/diseased boundary per host rather than across hosts.

For few-shot calibration, five support images per class ($K=5$) are drawn at random, 55 in total. The remaining 9,330 images serve as unseen queries. Neither the support nor the query images are used to pretrain the frozen backbones.

\subsubsection{PlantDoc In-The-Wild Benchmark (Field Regime)}
To test robustness under field conditions (solar glare, canopy shading, soil clutter, overlapping foliage, motion blur), we use the in-the-wild PlantDoc benchmark \cite{singh2020plantdoc}. It contains 2,588 foliar images from 13 botanical species across 27 pathology classes, split officially into a training repository ($\mathcal{D}_{\text{train}}^{\text{PlantDoc}}$, 2,336 images, all 27 classes) and a test benchmark ($\mathcal{D}_{\text{test}}^{\text{PlantDoc}}$, 252 handheld field images from Indian orchards under ambient light). Support sets come only from $\mathcal{D}_{\text{train}}^{\text{PlantDoc}}$ and query sets only from $\mathcal{D}_{\text{test}}^{\text{PlantDoc}}$, so $\mathcal{S} \cap \mathcal{Q} = \emptyset$. We evaluate in two regimes:
\begin{itemize}
    \item \textbf{5-way 5-shot episodic:} Following standard few-shot protocols \cite{snell2017prototypical,zhou2026vsf}, we run 50 independent Monte Carlo episodes. In each episode $e$, 5 classes are drawn without replacement ($\mathcal{C}_e \subset \{1, \dots, 27\}$), and the support set $\mathcal{S}_e$ has 5 labeled leaves per class (25 total) from $\mathcal{D}_{\text{train}}^{\text{PlantDoc}}$. The query set $\mathcal{Q}_e$ is every test image of those 5 classes (46.7 on average), with $\mathcal{S}_e \cap \mathcal{Q}_e = \emptyset$.
    \item \textbf{27-way all-class stress test:} All 27 classes are evaluated at once (chance: $1/27 \approx 3.70\%$). The model gets 5 support leaves per class (135 total) from $\mathcal{D}_{\text{train}}^{\text{PlantDoc}}$, expanded by 8-view augmentation to $N = 1{,}080$ feature vectors, and scores all 252 field queries in closed form ($\mathcal{S} \cap \mathcal{Q} = \emptyset$).
\end{itemize}

\subsection{Evaluation Protocols and Epidemiological Safety Metrics}
\label{sec:metrics}
Experiments on PlantVillage are evaluated over 30 independent Monte Carlo episodes with randomized support splits, reporting mean accuracy $\bar{x}$, standard deviation $\sigma$, 95\% confidence intervals (CI; $95\%\,\text{CI} = \bar{x} \pm 1.96 \cdot \sigma / \sqrt{30}$), and peak single-episode accuracy. Because misclassification costs in precision agriculture are asymmetric, we evaluate epidemiological safety metrics alongside standard accuracy:
\begin{itemize}
    \item \textbf{Top-1 Accuracy \& Macro F1-Score:} The F1-score is the harmonic mean of precision and recall, macro-averaged across all classes to prevent frequent classes (e.g., Potato Early Blight with 995 queries) from masking errors on small or healthy classes (e.g., Potato Healthy with 147 queries).
    \item \textbf{False-Negative Disease Escape Rate (FNR):} The proportion of diseased leaves erroneously diagnosed as healthy, where false-negative (FN) and true-positive (TP) counts define the rate: $\text{FNR} = \text{FN} / (\text{TP} + \text{FN})$. Undetected infections can lead to rapid pathogen spread in the field.
    \item \textbf{Cross-Species Confusion Rate:} The proportion of queries misattributed to an entirely different botanical family (e.g., classifying a potato leaf as an apple disease), indicating cross-taxonomic error.
\end{itemize}

To ensure unbiased cross-domain evaluation, the shared constants ($\gamma = 20.0$, $w_{\text{lda}} = 0.75$, and $\tau_{\text{conf}} = 0.35$) were calibrated strictly on held-out PlantVillage development episodes (Section~\ref{sec:calibration}) and carried to PlantDoc without domain-specific re-tuning. The anchor buffer is the one exception, and it is structural rather than tuned: PlantDoc yields only 252 test queries against 27 classes, so the 0.70 confidence threshold and the 20-anchor cap starve the buffer (Section~\ref{sec:transductive_validation} reports the same starvation at $\tau \ge 0.8$ in-domain). For the field protocol the buffer therefore admits every confident prediction and caps at two anchors per class. No accuracy-tuned constant differs between the two datasets. In transductive evaluation the main results are computed over a single query batch; Section~\ref{sec:ablation} additionally simulates sequential arrival. We report both static inductive and transductive results across all benchmarks to decouple baseline feature discrimination from dynamic adaptation gains.

\subsection{Edge Hardware Testbed}
\label{sec:edge_hardware_testbed}
On-device profiling is conducted on a Samsung Galaxy S24~FE handset (Exynos 2400e SoC, ARM64) running Android~17. Neural network inference executes strictly on the CPU via ONNX Runtime with per-channel INT8 quantization (4 intra-op threads; GPU and NNAPI delegates disabled), while closed-form linear algebra runs via the Efficient Java Matrix Library (EJML). Benchmark latencies and process memory (PSS) are averaged over 100 iterations following 10 warm-up runs, with thermal stability monitored across sustained evaluation. Detailed hardware specifications and component-level latency breakdowns are provided in Section~\ref{sec:edge_hardware} (Tables~\ref{tab:device_spec}--\ref{tab:adaptation_breakdown}).

